\section{Introduction}

Let $X$ be a smooth projective geometrically connected curve over
$\F_q$, let $F=\F_q(X)$, and let $G/F$ be a split connected adjoint
absolutely simple group. Fix a split Borel subgroup $B=TN$, and write
$\A_F$ for the adeles of $F$. A cuspidal automorphic representation
$\pi$ of $G(\A_F)$ is \emph{globally generic} if some vector has a
nonzero Whittaker coefficient
\[
 \int_{N(F)\backslash N(\A_F)} f(ng)\psi(n)^{-1}\,dn
\]
for a character $\psi$ nontrivial in every simple-root coordinate.
The generalized Ramanujan conjecture predicts tempered local
components for such representations. Here temperedness means that
the unitary local representation is weakly contained in the regular
representation. This is the generic temperedness prediction within the
broader framework of Arthur parameters \cite{Arthur1989,Shahidi2011}.

We prove the exceptional-group case, including the places where the
representation is ramified. The global input is the unramified
Ramanujan theorem of the companion paper \cite[Corollary~1.2]{UnramifiedRamanujan}.

\begin{theorem}\label{thm:main}
Let $G/F$ be split connected adjoint absolutely simple of type
$G_2$, $F_4$, $E_6$, $E_7$, or $E_8$. If
$\pi=\bigotimes'_v\pi_v$ is a complex globally generic cuspidal
automorphic representation of $G(\A_F)$, then $\pi_v$ is tempered
for every place $v$. There is no restriction on the characteristic
of $F$, the ramification depth, or the local representation type.
\end{theorem}

The distinction between spherical and ramified components is
substantial. An unramified parameter is determined by its Satake
class, whereas at a ramified place a Weil--Deligne parameter also
contains a nilpotent monodromy operator. The local-global theorem
available for general reductive groups identifies only the
semisimple Weil parameter. The argument below supplies the
additional monodromy comparison needed for temperedness.

\subsection{Context and the local problem}

Over function fields, Drinfeld proved Ramanujan for $\GL_2$
\cite{Dri88}, and Laurent Lafforgue proved it for $\GL_n$
\cite[Th\'eor\`eme~VI.10(i)]{Laf02}. For other reductive groups,
genericity selects the part of the cuspidal spectrum for which
temperedness is expected. The Langlands--Shahidi method relates this
condition to local coefficients and automorphic $L$-functions.
Lomel\'i developed this method over function fields and proved
generic Ramanujan for split classical groups and quasi-split unitary
groups \cite{Lom09,LomeliLS,Lom19}. In particular, the later split
classical results include characteristic $2$
\cite[Theorem~7.1(i)]{Lom19}.

Vincent Lafforgue's excursion operators attach semisimple global
parameters to cuspidal automorphic representations of general
reductive groups \cite{VLafforgue}. Genestier--Lafforgue construct
semisimple local parameters and prove local-global compatibility
at every place \cite{GenestierLafforgue}. Gan--Harris--Sawin, with
an appendix by Beuzart-Plessis, show that the local parameter of a
tempered representation admits an essentially tempered completion
\cite[Theorem~1.2 and Corollary~1.3]{GHS}. These results make it
possible to compare the global and representation-theoretic sources
of monodromy without assuming a full local correspondence.

On the unramified side, Sawin--Templier prove temperedness under a
monomial geometric supercuspidal hypothesis and a cyclic base-change
hypothesis \cite{SawinTemplier2021}. Ciubotaru--Harris obtain temperedness at
unramified places under hypotheses including a generic unramified
place and a tempered place \cite{CH26}. The companion result used
here gives the unramified conclusion for split adjoint absolutely
simple groups from a generic unramified component alone
\cite[Corollary~1.2]{UnramifiedRamanujan}. The present paper proves
the local implication that extends this conclusion to the places
of level.

Two established principles guide that implication. First, genericity
is related to regularity at $1$ of the adjoint $L$-factor, as formulated
in the Gross--Prasad/Rallis conjecture
\cite[Conjecture~2.6]{GrossPrasad1992}; see, for
example, the criterion proved by Gan--Ichino under local-correspondence
and local-factor hypotheses \cite[Appendix~B, Proposition~B.1]{GanIchino2016}.
Second, adjoint regularity is equivalent to openness of the
monodromy orbit over a fixed semisimple parameter
\cite[Proposition~3.5]{CDFZ2025} and \cite[Proposition~6.10]{DHKM2026}.
We give the short Lie-algebra proof of the implication needed here.
The broader principle that purity strongly constrains a completion
of a semisimple Weil representation also appears in
Taylor--Yoshida \cite[Section~1]{TaylorYoshida2007}.

Our task is to obtain the required adjoint regularity for arbitrary
generic inducing data using the semisimple correspondence alone.
We compare the \emph{total} local coefficient with a product of
Weil--Deligne $L$-factor ratios. Multiplicativity, rank-one Tate
factors, the globalization theorem of Gan--Lomel\'i, and the crude
functional equation of Lomel\'i provide the comparison
\cite{GanLomeli,LomeliLS}. An elementary cancellation shows that
these ratios do not depend on monodromy up to nonvanishing monomials.
This is exactly the amount of local-factor compatibility required
by the argument. Genestier--Lafforgue also construct local $\gamma$-factors
from semisimple parameters and characterize them through global
functional equations \cite[Section~7]{GenestierLafforgue}.

\subsection{The proof and its reusable local criterion}

The proof has a local part and a global source of purity. At a fixed
place, let $r$ be the semisimple Weil parameter of a generic local
representation. A Weil--Deligne completion is a nilpotent operator
$N$ satisfying $\Ad(r(w))N=|w|N$. The local result is
Theorem~\ref{thm:local-criterion}: if a completion has adjoint
representation pure of weight zero, then the representation is
tempered. Purity here prescribes Frobenius eigenvalue sizes along
monodromy chains, with weights symmetric about zero. This criterion applies to every split adjoint absolutely
simple group over a local function field.

To see the mechanism, suppose the representation is not tempered.
Its Langlands datum consists of a tempered representation $\sigma$
of a proper Levi subgroup and a strictly positive real exponent
$\nu$. A tempered completion for $\sigma$, twisted by $\nu$, gives
a completion $(r,N_M)$. The long intertwiner has generic image, so
its scalar on Whittaker coinvariants is nonzero. The local-coefficient
comparison then makes the adjoint $L$-factor of $(r,N_M)$ regular at
$1$. A completion $(r,N_*)$ with pure adjoint has the same regularity. Both
operators therefore lie in the unique open orbit of the centralizer
of $r$, and are conjugate. Figure~\ref{fig:comparison} displays this
comparison.

\begin{figure}[tbp]
\centering
\begin{tikzpicture}[
 box/.style={draw,rounded corners=2pt,align=center,text width=5.6cm,
             inner sep=7pt,font=\small},
 arr/.style={-{Stealth[length=2mm]},thick}]
 \node[box] (global) at (0,0)
  {Completion $(r,N_*)$ with pure adjoint\\
   Adjoint $L$-factor regular at $1$};
 \node[box] (local) at (6.9,0)
  {Generic quotient of\\positive Langlands data\\
   Completion $(r,N_M)$ with the same regularity};
 \node[box,text width=8.7cm] (orbit) at (3.45,-2.1)
  {A unique open orbit over the fixed Weil parameter $r$\\
   $N_*\sim N_M$, so adjoint purity passes to $(r,N_M)$};
 \draw[arr] (global.south) -- (orbit.north west);
 \draw[arr] (local.south) -- (orbit.north east);
\end{tikzpicture}
\caption{The local comparison uses two independently obtained
monodromy operators. Adjoint regularity identifies their orbits
over the fixed Weil parameter $r$.}
\label{fig:comparison}
\end{figure}

Purity now passes to the positive dual nilradical, a direct summand
of the adjoint Weil--Deligne representation. A pure weight-zero
representation has Frobenius determinant of absolute value $1$.
The positive exponent $\nu$ makes that same determinant strictly
smaller than $1$. This contradiction proves the local criterion.

Globally, the companion theorem makes every unramified component
tempered. The adjoint local system attached to an occurring
excursion parameter is consequently pointwise pure of weight zero.
Deligne's theorem on local monodromy for curves
\cite[Th\'eor\`eme~1.8.4]{DeligneWeilII} supplies a pure adjoint
completion at each place. Semisimple local-global compatibility
identifies its Weil part with the parameter of $\pi_v$, and the
local criterion applies.

Section~\ref{sec:parameters} fixes conventions and states the
parameter inputs. Section~\ref{sec:monodromy} proves the monodromy
and purity lemmas. Section~\ref{sec:coefficients} establishes the
local-coefficient comparison, including its behavior under
specialization. Section~\ref{sec:local-criterion} proves the local
criterion, and Section~\ref{sec:global} applies it to the global
representation.
